\documentclass[11pt,a4paper]{amsart}
\usepackage{cgeometry}
\addbibresource{ComplexGeometry.bib}
\renewcommand*{\bibfont}{\footnotesize}

\allowdisplaybreaks

\newcommand{\dHaus}{d_{\mathrm{Haus}}}
\newcommand{\dS}{d_{S}}
\newcommand{\Kconv}{\mathcal K_n}
\newcommand{\MultiProj}{\operatorname{MultiProj}}
\newcommand{\inival}{\operatorname{in}_{\nu}}
\DeclareMathOperator{\wt}{wt}

\title{Mixed volumes of Okounkov bodies}
\author{Mingchen Xia}


\begin{document}
\begin{abstract}
Let $X$ be a connected smooth projective variety of dimension $n$.
We prove that the movable product of any big line bundles $L_1,\ldots,L_n$ equals $n!$ times the supremum of the mixed volumes of their Okounkov bodies over all valuations.
\end{abstract}
\maketitle

\section{Introduction}

Let $X$ be a connected smooth projective variety of dimension $n$.  For a
big line bundle $L$ and a surjective valuation
\[
 \nu\colon \mathbb C(X)^\times\longrightarrow\mathbb Z^n_{\lex}
\]
that is trivial on $\mathbb C^\times$, write
$\Delta_\nu(L)\subseteq\mathbb R^n$ for the associated Okounkov body.  Its
Euclidean volume recovers the volume, or equivalently the top movable
self-intersection, of $L$:
\begin{equation}\label{eq:diagonal-volume}
 n!\vol \Delta_\nu(L) =\vol(L)=\langle L^n\rangle.
\end{equation}
For the full valuation class used here, Okounkov bodies depend only on numerical classes \cite[Proposition~10.3.2]{Xiabook}; the flag-valuation case is \cite[Proposition~4.1]{LM09}.  Conversely, Jow proved that two big divisors
on a normal complex projective variety are numerically equivalent if their Okounkov bodies agree for every admissible flag
\cite[Theorem~A]{Jow10}.  Thus the collection of all Okounkov bodies is a faithful numerical invariant.

The first conclusion of this paper is that this collection also recovers mixed movable products. 
\begin{theorem}\label{thm:main1}
    Let $L_1,\ldots,L_n$ be big line bundles on $X$, then
    \begin{equation}\label{eq:movable-main}
        \langle L_1,\ldots,L_n\rangle =n!\sup_\nu\vol\Bigl(\Delta_\nu(L_1),\ldots,\Delta_\nu(L_n)\Bigr),
    \end{equation}
    where $\nu$ ranges over the valuations above.
\end{theorem}
Here $\vol$ denotes mixed volume, normalized by $\vol(K,\ldots,K)=\vol(K)$. For ample line bundles, Wilms proved that a suitable single admissible flag realizes the mixed intersection number \cite[Theorem~1.1]{Wil25}.

We now pass to its singular-metric refinement.  A Hermitian big line bundle $(L,h)$ will mean a big line bundle endowed with a singular psh metric with positive volume (in the sense of Cao). In \cite{Xia21}, the author introduced the partial Okounkov body
$\Delta_\nu(L,h)\subseteq\mathbb R^n$ extending the usual Okounkov body construction. The volume theorem \cite[Theorem~A]{Xia21} gives
\begin{equation}\label{eq:partial-diagonal-volume}
 n!\vol \Delta_\nu(L,h) =\vol \ddc h,
\end{equation}
where the volume on the right-hand side is taken in the sense of Cao \cite{Cao14}.
Moreover, the family of partial Okounkov bodies over all rank-$n$ valuations determines the $\mathcal{I}$-singularity type of $h$
\cite[Theorem~B]{Xia21}, a metric analogue of Jow's theorem.

A formula of this type was first proposed in \cite[Conjecture~8.8]{Xia21}. We prove it here.
\begin{theorem}\label{thm:main}
Let $X$ be a connected smooth projective variety of dimension $n$.  For $1\leq i\leq n$, let $L_i$ be a big line bundle and let $h_i$ be a psh metric on $L_i$ whose volume is positive. Then
\begin{equation}\label{eq:main}
 \vol\left(\ddc h_1,\ldots,\ddc h_n\right)=n!\sup_\nu\vol\Bigl(\Delta_\nu(L_1,h_1),\ldots,\Delta_\nu(L_n,h_n)\Bigr).
\end{equation}
Here $\nu$ runs over all valuations as above.
\end{theorem}
Again, the mixed volume is taken in the sense of Cao; by \cite[Definition~3.2, Proposition~3.3, and Theorem~3.5]{Xiabdiv}, it coincides with the mixed volume in the sense of Darvas--Xia recalled in \cref{sec:preliminaries}. Taking each $h_i$ to have minimal singularities reduces \cref{thm:main} to \cref{thm:main1}, while its diagonal specialization is \eqref{eq:partial-diagonal-volume}.

The proof of the $\leq$ direction in \eqref{eq:main} can be easily reduced to Wilms' theorem \cite[Theorem~1.1]{Wil25}. Conversely, the $\geq$ direction can be reduced to the $\geq$ direction in \eqref{eq:movable-main}. We shall apply a suitable toric degeneration to further reduce to the toric case, where the results are essentially known.

\subsection*{Acknowledgments}
The author used the \href{https://github.com/frenzymath/Rethlas}{Rethlas agent}, backed by an OpenAI GPT-5.6 Sol model, for proof exploration and consistency checks during the development of this work.  AI-generated suggestions were treated as
provisional; the author assumes full responsibility for every argument and for the final text.

The author is supported by the National Key R\&D Program of China under Grant No.~2025YFA1018200.

\section{Preliminaries}\label{sec:preliminaries}

We follow the language and notation of \cite{Xiabook}.  All products of closed positive currents are non-pluripolar products.  We use
$\ddc=(2\pi)^{-1}\ii\partial\bar\partial$, and $\mathbb N=\{0,1,2,\ldots\}$. The notation $\mathbb Z^n_{\lex}$ means $\mathbb Z^n$ with the lexicographic order, and $e_1,\ldots,e_n$ denotes its standard basis. A superscript ${}^\times$ means that zero is omitted, and $\Conv A$ denotes the convex hull of a subset $A\subseteq\mathbb R^n$.  We use $\mathbb P(V)$ for the space of lines in a vector space $V$.

\subsection{Singularities and partial Okounkov bodies}

Fix a smooth Hermitian metric $h^0$ on a big line bundle $L$ and write
\[
 \theta\coloneqq c_1(L,h^0),\qquad h=h^0\mathrm e^{-\varphi},\qquad \theta_\varphi\coloneqq \theta+\ddc\varphi,
\]
where $\PSH(X,\theta)$ is the set $\theta$-psh functions, and
\[
 \PSH(X,\theta)_{>0} \coloneqq \left\{u\in\PSH(X,\theta):\int_X\theta_u^n>0\right\}.
\]
For $u,v\in\PSH(X,\theta)$, write $u\preceq v$ if
$u\leq v+O(1)$.  The model envelope is
\[
 P_\theta[u]
 \coloneqq \sup{}^*\{v\in\PSH(X,\theta):v\leq0,\ v\preceq u\},
\]
where $\sup{}^*$ denotes the upper-semicontinuous regularization of the pointwise supremum.  A positive-mass potential $u$ is a model potential if $P_\theta[u]=u$.  Two positive-mass potentials are $P$-equivalent if their model envelopes agree.  They are $\mathcal{I}$-equivalent if $\mathcal{I}(\lambda u)=\mathcal{I}(\lambda v)$ for every $\lambda>0$, where $\mathcal{I}(\lambda u)$ is the multiplier ideal sheaf. The $\mathcal{I}$-model envelope is
\[
 P_\theta[u]_{\mathcal{I}} \coloneqq \sup{}^*\{v\in\PSH(X,\theta):v\leq 0,\ v\text{ is }\mathcal{I}\text{-equivalent to }u\}.
\]
For an arbitrary $\varphi\in\PSH(X,\theta)$, its volume in the sense of Cao is
\[
    \vol \theta_{\varphi} \coloneqq \int_X\theta_{P_\theta[\varphi]_{\mathcal{I}}}^n.
\]
A potential $\varphi\in \PSH(X,\theta)_{>0}$ is called $\mathcal{I}$-good if
\[
 P_\theta[\varphi]=P_\theta[\varphi]_{\mathcal{I}}.
\]
More generally, if $L_1,\ldots,L_n$ are big line bundles with $\theta_i$ being a closed smooth real $(1,1)$-form in $c_1(L_i)$ and $\varphi_i\in \PSH(X,\theta_i)$ satisfying $\vol(\theta_i+\ddc\varphi_i)>0$, we write
\[
\vol\left( \theta_1+\ddc \varphi_1,\ldots, \theta_n+\ddc \varphi_n\right)=\int_X \left(\theta_1+\ddc P_{\theta_1}\left[\varphi_1\right]_{\mathcal{I}}\right)\wedge \cdots \wedge \left(\theta_n+\ddc P_{\theta_n}\left[\varphi_n\right]_{\mathcal{I}}\right).
\]
The mixed volume defined here coincides with that defined by Cao
\cite[Definition~3.2, Proposition~3.3, and Theorem~3.5]{Xiabdiv}.


By \cite[Theorem~7.1.1]{Xiabook}, $\mathcal{I}$-goodness is equivalent both to the mass--volume identity
\[
 \int_X\theta_\varphi^n=\vol \theta_{\varphi}
\]
and to the existence of potentials $\varphi_j$ with analytic singularities such that
\[
 \varphi_j\xrightarrow{\dS}\varphi.
\]
Here $\dS$ is the pseudometric on singularity types introduced in \cite{DDNLmetric}, see also \cite[Definition~6.2.1]{Xiabook}. The approximants may be chosen so that every $\theta_{\varphi_j}$ is a K\"ahler current.  Following \cite[Definition~1.6.1]{Xiabook}, a potential has analytic singularities if locally it is of the form
\[
 c\log\bigl(|f_1|^2+\cdots+|f_N|^2\bigr)+O(1), \qquad c\in\mathbb Q_{>0},
\]
where the $f_j$ are holomorphic and $O(1)$ is locally bounded.

Let $\nu\colon\mathbb C(X)^\times\to\mathbb Z^n_{\lex}$ be surjective and
trivial on $\mathbb C^\times$.  Its center exists on $X$ by properness.
Choose a rational section $\tau$ of $L$ which is regular and nowhere
vanishing near the center.  For a section $s$ of $L^k$, set
$\nu_\tau(s)\coloneqq \nu(s/\tau^k)$, and suppress the subscript $\tau$ below.
For every $k\geq1$ for which the space below is non-zero, set
\begin{equation}\label{eq:level-body}
 \Delta_{\nu,k}(\theta,\varphi)
 \coloneqq \Conv\left\{k^{-1}\nu(s): s\in \mathrm{H}^0\bigl(X,L^k\otimes\mathcal{I}(k\varphi)\bigr)^{\times}\right\}.
\end{equation}
Let $\Gamma(\theta,\varphi)\subseteq\mathbb Z^n\times\mathbb N$ consist of
$(0,0)$ and all pairs $(\nu(s),k)$ occurring above.  By
\cite[Theorem~10.3.1]{Xiabook}, $\Gamma(\theta,\varphi)$ is an almost
semigroup.  Using the continuous Okounkov-body map $\Delta$ of
\cite[Theorem~10.1.2]{Xiabook}, define
$\Delta_\nu(\theta,\varphi)=\Delta_\nu(L,h)
\coloneqq \Delta(\Gamma(\theta,\varphi))$.  Changing the rational
trivialization translates all the bodies in \eqref{eq:level-body}; none of
the mixed volumes below is affected.

We record the properties used in the proof.  We write
$\varphi\preceq_{\mathcal{I}}\psi$ when
$\mathcal{I}(\lambda\varphi)\subseteq\mathcal{I}(\lambda\psi)$ for every $\lambda>0$.
Let $\Kconv$ be the space of non-empty compact convex subsets of
$\mathbb R^n$, and let $\dHaus$ be the Hausdorff distance on $\Kconv$.
We write $\vol(K)$ for the Euclidean $n$-volume of
$K\in\Kconv$.
First,
\begin{equation}\label{eq:partial-body-volume}
 \vol \Delta_\nu(\theta,\varphi) =\frac{1}{n!}\vol \theta_{\varphi}
\end{equation}
by \cite[Corollary~10.3.2]{Xiabook}.  The right-hand side is independent of
$\nu$.  Second, the map
\[
 \Delta_\nu(\theta,\cdot)\colon
 (\PSH(X,\theta)_{>0},\dS)\longrightarrow(\Kconv,\dHaus)
\]
is continuous by \cite[Theorem~10.3.2]{Xiabook}.  Third, for another Hermitian
big line bundle $(L',h')$, singularity comparison and tensor products give
\begin{align}
 \varphi\preceq_{\mathcal{I}}\psi
 &\Longrightarrow
 \Delta_\nu(\theta,\varphi)\subseteq
 \Delta_\nu(\theta,\psi),                                      \label{eq:body-monotonicity}\\
 \Delta_\nu(L,h)+\Delta_\nu(L',h')
 &\subseteq
 \Delta_\nu(L\otimes L',h\otimes h'),                         \label{eq:body-Minkowski}
\end{align}
by \cite[Propositions~10.3.8 and~10.3.10]{Xiabook}.  In particular, the
partial body is contained in the ordinary Okounkov body $\Delta_\nu(L)$,
obtained from the same construction using the complete linear series
$\mathrm{H}^0(X,L^k)$.
The latter satisfies
\begin{equation}\label{eq:ordinary-volume}
 \vol \Delta_\nu(L)=\frac{1}{n!}\vol(L)
\end{equation}
by \cite[Corollary~10.3.1]{Xiabook}, where $\vol(L)\coloneqq \lim_{k\to\infty}n!k^{-n}h^0(X,L^k)$.

We will also use basic facts about model envelopes.  By
\cite[Corollary~3.1.2]{Xiabook}, $P_\theta[\varphi]$ is a model potential.
The equality $\dS(\varphi,P_\theta[\varphi])=0$ follows from
\cite[Proposition~6.2.2]{Xiabook}, and $P$-equivalent potentials are
$\mathcal{I}$-equivalent by \cite[Proposition~3.2.9]{Xiabook}.  Consequently, replacing
a potential by its model envelope changes neither its $\dS$-class nor its
partial Okounkov body.  If $\varphi$ has analytic singularities, then so does
$P_\theta[\varphi]$ by \cite[Proposition~3.2.10]{Xiabook}.
We will separately use the $\mathcal I$-model replacement.  By
\cite[Proposition~3.2.6]{Xiabook}, $\varphi$ and
$P_\theta[\varphi]_{\mathcal I}$ are $\mathcal I$-equivalent, so they have
the same partial Okounkov bodies.  If $\vol\theta_\varphi>0$, then
\[
 \int_X\theta_{P_\theta[\varphi]_{\mathcal I}}^n
 =\vol\theta_\varphi>0,
\]
and $P_\theta[\varphi]_{\mathcal I}$ is an $\mathcal I$-good model
potential by \cite[Example~7.1.2]{Xiabook}.

\begin{lemma}\label{lem:potential-sandwich}
Let $u_j,u\in\PSH(X,\theta)_{>0}$ be model potentials such
that $u_j\xrightarrow{\dS}u$. The sequence $(u_j)_j$ has a further subsequence, still indexed by $j$, for which there are increasing model potentials $\eta_{j}$ and decreasing model potentials $\psi_{j}$ satisfying
\[
 \eta_{j}\leq u_{j}\leq\psi_{j},\qquad
 \eta_{j}\leq u\leq\psi_{j},\qquad
 \eta_{j}\xrightarrow{\dS}u,\qquad
 \psi_{j}\xrightarrow{\dS}u.
\]
\end{lemma}

\begin{proof}
By \cite[Corollary~6.2.5]{Xiabook}, the masses $\int_X\theta_{u_j}^n$ have a positive lower bound after finitely many terms. So we can apply \cite[Proposition~6.2.3]{Xiabook} and then take the $P$-envelopes to obtain $\eta_j$ and $\psi_j$.  It remains to compare them with the limit.  Fix $j$.  For every $k\geq j$,
\[
 \eta_j\leq\eta_k\leq u_k.
\]
Apply the closure of the $P$-order under $\dS$-convergence
\cite[Proposition~6.2.5]{Xiabook} to the constant sequence $\eta_j$ and
the tail $(u_k)_{k\geq j}$.  Their model envelopes satisfy
$P_\theta[\eta_j]\leq P_\theta[u]$.  Since both potentials are model, this
is $\eta_j\leq u$.  Likewise, $u_k\leq\psi_k\leq\psi_j$ for $k\geq j$;
order closure applied to the tail $(u_k)_{k\geq j}$ and the constant
sequence $\psi_j$ gives $P_\theta[u]\leq P_\theta[\psi_j]$, hence
$u\leq\psi_j$.  Together with $\eta_j\leq u_j\leq\psi_j$, this proves the
assertion.
\end{proof}

\subsection{Mixed volumes and movable intersections}

We retain the mixed-volume normalization fixed in \cref{thm:main}.  Mixed
volume is symmetric, multilinear for Minkowski addition, monotone,
translation invariant, and continuous for the Hausdorff topology; see
\cite[Section~5.1]{Sch14}.

For big numerical divisor classes $\alpha_1,\ldots,\alpha_n$ on a smooth projective $X$, we write
$\langle \alpha_1,\ldots,\alpha_n \rangle$ for their movable intersection number.  The nef-approximation description is \cite[Definition~2.5 and Proposition~2.13]{BFJ09}.  If $\mu\colon Y\to X$ is a modification with $Y$ smooth and
\[
 \mu^*\alpha_i=\beta_i+\{E_i\},
 \qquad \beta_i\ \text{nef},\quad E_i\textup{ is an effective }\mathbb{R}\textup{-divisor},
\]
then
\begin{equation}\label{eq:movable-admissible}
 \beta_1\cdots \beta_n\leq\langle \alpha_1,\ldots,\alpha_n\rangle.
\end{equation}
For nef $\alpha_i$, the movable and ordinary intersection numbers agree by
\cite[Proposition~2.12]{BFJ09}.

\section{An algebraic upper bound}

\begin{lemma}\label{lem:finite-level}
Let $F/\mathbb C$ be a finitely generated field of transcendence degree $n$,
let $X$ be an integral projective model with $\mathbb C(X)=F$, and let $\nu$
be a surjective valuation which is trivial on $\mathbb C^\times$, written as
\[
 \nu\colon F^\times\longrightarrow\mathbb Z^n_{\lex}.
\]
Let $W_1,\ldots,W_n\subset F$ be non-zero finite-dimensional vector spaces.
Evaluation of rational functions defines
\[
 \Phi_i\colon X\dashrightarrow\mathbb P(W_i^\vee),\qquad
 x\longmapsto[\operatorname{ev}_{i,x}],\qquad
 \operatorname{ev}_{i,x}(f)\coloneqq f(x).
\]
This map is defined wherever all elements of $W_i$ are regular and
$\operatorname{ev}_{i,x}\neq0$.  Put
\[
 \Phi\coloneqq (\Phi_1,\ldots,\Phi_n)\colon
 X\dashrightarrow\prod_{i=1}^n\mathbb P(W_i^\vee),
\]
and let $Y$ be the Zariski closure of its image.
Assume that $\Phi$ is birational onto $Y$ and that
\[
 \Lambda\coloneqq \operatorname{span}_{\mathbb Z}
 \{\nu(f)-\nu(g):f,g\in W_i^\times,\ 1\leq i\leq n\}=\mathbb Z^n.
\]
Put $P_i\coloneqq \Conv\nu(W_i^\times)$.  If $\pi_i$ is the projection of the
ambient product onto $\mathbb P(W_i^\vee)$, put
\[
 H_i\coloneqq c_1\bigl(\pi_i^*\mathcal O_{\mathbb P(W_i^\vee)}(1)\bigr).
\]
 Then
\begin{equation}\label{eq:finite-level}
 n!\vol(P_1,\ldots,P_n)\leq(H_1\cdots H_n\cdot Y).
\end{equation}
\end{lemma}

\begin{proof}
If $\dim_{\mathbb C}W_i=1$ for some $i$, then $P_i$ is a point and
$H_i=0$, since the $i$-th factor is $\mathbb P^0$.  Both sides of
\eqref{eq:finite-level} are therefore zero.  We may henceforth assume that
$\dim_{\mathbb C}W_i\geq2$ for every $i$; in particular, the general
hyperplanes used below exist in every factor.

For $a\in\mathbb Z^n$, let
\[
 \mathcal F_{\succeq a}\coloneqq \{0\}\cup\{f\in F^\times:\nu(f)\succeq a\},
 \qquad
 \mathcal F_{\succ a}\coloneqq \{0\}\cup\{f\in F^\times:\nu(f)\succ a\},
\]
where $\succeq$ is the lexicographic order.  The associated graded ring and
the initial form of $f\in F^\times$ are
\[
 \operatorname{gr}_\nu F
 \coloneqq \bigoplus_{a\in\mathbb Z^n}
    \mathcal F_{\succeq a}/\mathcal F_{\succ a},
 \qquad
 \inival(f)\coloneqq f\bmod\mathcal F_{\succ\nu(f)}.
\]
The quotient $\mathcal F_{\succeq a}/\mathcal F_{\succ a}$ is called the
leaf of $\nu$ at $a$.
Since $\nu$ has rational rank $n$, Abhyankar's inequality forces its residue
field to be $\mathbb C$; hence every non-zero graded piece is one-dimensional
over $\mathbb C$.  Indeed, if two non-zero elements have the same value,
the residue of their quotient is a scalar, and subtracting the corresponding
scalar multiple raises the value.  Equivalently, every non-zero finite-dimensional subspace
$V\subset F$ satisfies
\[
 \#\nu(V^\times)=\dim_{\mathbb C}V;
\]
compare \cite[Lemma~1.3]{LM09} and \cite[Section~10.3]{Xiabook}.

Put $r_i\coloneqq \dim_{\mathbb C}W_i-1$.  Choose one
$f_{ij}\in W_i^\times$ above each value in $\nu(W_i^\times)$ and set
\[
 a_{ij}\coloneqq \nu(f_{ij}),\qquad 0\leq j\leq r_i.
\]
The $a_{ij}$ are pairwise distinct, so the valuation axiom makes the
$f_{ij}$ linearly independent.  They therefore form a basis of $W_i$, and
\[
 \nu(W_i^\times)=\{a_{i0},\ldots,a_{ir_i}\}.
\]

\smallskip
\noindent\emph{The initial degeneration.}
Choose $u_\ell\in F^\times$ with $\nu(u_\ell)=e_\ell$.  One-dimensionality
of the leaves gives a non-canonical graded isomorphism.  More precisely, for
$a=(a_1,\ldots,a_n)\in\mathbb Z^n$, write
$u^a\coloneqq u_1^{a_1}\cdots u_n^{a_n}$, and write
\[
 \mathbb C[\mathbb Z^n]
 \coloneqq \bigoplus_{a\in\mathbb Z^n}\mathbb C\chi^a,
 \qquad \chi^a\chi^{a'}=\chi^{a+a'},
\]
for the group algebra of $\mathbb Z^n$.  Then
\[
 \mathbb C[\mathbb Z^n]\xrightarrow{\sim}\operatorname{gr}_\nu F,
 \qquad
 \chi^a\longmapsto\inival(u^a).
\]
After rescaling the $f_{ij}$, which merely changes the homogeneous
coordinates, we may and do assume that
\[
 \inival(f_{ij})=\chi^{a_{ij}}.
\]

Let $x_{ij}$ be the homogeneous coordinate on $\mathbb P(W_i^\vee)$ dual to
$f_{ij}$.  In these coordinates,
\[
 \Phi_i(x)=[f_{i0}(x):\cdots:f_{ir_i}(x)].
\]

Let $\epsilon_1,\ldots,\epsilon_n$ be the standard basis of the multidegree lattice $\mathbb Z^n$.  The multihomogeneous coordinate ring of
the ambient product is
\[
 S\coloneqq \mathbb C[x_{ij}:1\leq i\leq n,\ 0\leq j\leq r_i], \qquad \deg x_{ij}\coloneqq \epsilon_i.
\]
Let $I\subseteq S$ be the multihomogeneous prime ideal of $Y$.  Equivalently, a multihomogeneous polynomial $G$ belongs to $I$ exactly when $G(f_{ij})=0$ in $F$.

Recall that each $a_{ij}$ is a column vector in $\mathbb Z^n$.  Put $N\coloneqq \sum_{i=1}^n(r_i+1)$, and define $\mathsf A$ to be the $n\times N$ matrix whose column indexed by $(i,j)$ is $a_{ij}$.  Equivalently,
\[
 \mathsf A\coloneqq \bigl[a_{10}\ \cdots\ a_{1r_1}\mid\cdots\mid
                   a_{n0}\ \cdots\ a_{nr_n}\bigr]
 \in\mathbb Z^{n\times N}
\,.
\]
For
$\alpha=(\alpha_{ij})\in\mathbb N^N$, write
$x^\alpha\coloneqq \prod_{i,j}x_{ij}^{\alpha_{ij}}$, and define
\[
 \wt_{\mathsf A}(x^\alpha)
 \coloneqq \sum_{i,j}\alpha_{ij}a_{ij}\in\mathbb Z^n.
\]
These weights are compared lexicographically.  If
$0\neq G=\sum_\alpha\gamma_\alpha x^\alpha$, let
$\operatorname{in}_{\mathsf A}(G)$ be the sum of the terms of least
$\mathsf A$-weight, and
put
\[
 Q\coloneqq \operatorname{in}_{\mathsf A}(I)
 \coloneqq \bigl(\operatorname{in}_{\mathsf A}(G):0\neq G\in I\bigr).
\]
Let $t_1,\ldots,t_n$ be algebraically independent variables of
multidegrees $\deg t_i=\epsilon_i$, and define
\[
 J\coloneqq \ker\left(S\longrightarrow
 \mathbb C[\mathbb Z^n][t_1,\ldots,t_n],\quad
 x_{ij}\longmapsto\chi^{a_{ij}}t_i\right).
\]
If $G\in I$ is multihomogeneous, then $G(f_{ij})=0$.  Taking initial forms
of the terms of least valuation in this equality gives
$\operatorname{in}_{\mathsf A}(G)\in J$.  Since $I$ is multihomogeneous, the initial
forms of its multihomogeneous elements generate $Q$; hence
\[
 Q\subseteq J.
\]
This is the argument of \cite[Lemma~2.16]{KM19}.

By \cite[Proposition~8.10]{KM19}, there is a rational weight vector
$\omega=(\omega_{ij})$ such that
$Q=\operatorname{in}_\omega(I)$, where $\operatorname{in}_\omega$ is
defined by taking terms of least $\omega$-weight.  After clearing denominators,
the $\omega_{ij}$ may be taken integral.  The weight filtration on $S/I$
respects the $\mathbb N^n$-grading and satisfies
\[
 \operatorname{gr}_\omega(S/I)
 \simeq S/\operatorname{in}_\omega(I)=S/Q.
\]
Each multihomogeneous piece is finite-dimensional, so passing to its
associated graded preserves dimension.  Hence
\begin{equation}\label{eq:hilbert-functions}
 \dim_{\mathbb C}(S/I)_{\mathbf d}
 =\dim_{\mathbb C}(S/Q)_{\mathbf d}
 \qquad(\mathbf d\in\mathbb N^n).
\end{equation}

\smallskip
\noindent\emph{The toric component.}
Let $T\coloneqq(\mathbb C^\times)^n$.  For
$z=(z_1,\ldots,z_n)\in T$ and
$a=(a_1,\ldots,a_n)\in\mathbb Z^n$, write
$z^a\coloneqq z_1^{a_1}\cdots z_n^{a_n}$, and define
\[
 \tau\colon T\longrightarrow\prod_{i=1}^n\mathbb P(W_i^\vee),\qquad
 z\longmapsto
 \left([z^{a_{i0}}:\cdots:z^{a_{ir_i}}]\right)_{i=1}^n.
\]
If $G\in S$ is multihomogeneous of multidegree
$\mathbf d=(d_1,\ldots,d_n)$, then its image under the homomorphism defining
$J$ is
\[
 t_1^{d_1}\cdots t_n^{d_n}G(\chi^{a_{ij}}).
\]
Thus $J$ is precisely the multihomogeneous ideal of the closure of
$\tau(T)$, and
\[
 Z\coloneqq\MultiProj(S/J)=\overline{\tau(T)}.
\]
On the dense torus of the ambient product, $\tau$ is defined by the
characters $z^{a_{ij}-a_{i0}}$.  Their exponents generate
$\Lambda=\mathbb Z^n$, so the induced homomorphism of character rings is
surjective.  Hence $\tau$ is a closed immersion into that torus.  We identify
$T$ with its image, a dense open subset of $Z$; in particular, $Z$ is
irreducible of dimension $n$.

Put $Y_0\coloneqq\MultiProj(S/Q)$.  The inclusion $Q\subseteq J$ gives a
closed immersion $Z\hookrightarrow Y_0$.  The ideal $Q$ need not be
saturated with respect to the multigraded irrelevant ideal, but this causes
no change in the argument: passing to $\MultiProj$ removes precisely the
irrelevant torsion, and the multigraded Serre theorem identifies the
sufficiently positive graded pieces with those of the corresponding
sheafification.  Thus \eqref{eq:hilbert-functions} gives a common
multigraded Hilbert polynomial for $Y$ and $Y_0$.  Its restriction to the
diagonal is the Hilbert polynomial of the ample line bundle
$\mathcal O(1,\ldots,1)$; hence
$\dim Y_0=\dim Y=n$.  Consequently, $Z$ is a top-dimensional irreducible
component of $Y_0$.

For an $n$-dimensional closed subscheme $V$ of the ambient product, the part
of total degree $n$ of its multigraded Hilbert polynomial is
\[
 \frac{1}{n!}
 \left(\left(\sum_{i=1}^n d_iH_i\right)^n\cdot[V]_n\right),
\]
where $[V]_n$ is its $n$-dimensional fundamental cycle.  Comparing the
coefficient of $d_1\cdots d_n$ yields
\[
 (H_1\cdots H_n\cdot Y)
 =(H_1\cdots H_n\cdot[Y_0]_n).
\]
Write
\[
 [Y_0]_n=\sum_Wm_W[W],\qquad
 m_W\coloneqq\ell\bigl(\mathcal O_{Y_0,\eta_W}\bigr)\geq1.
\]
Here $W$ runs over the $n$-dimensional irreducible components of $Y_0$ and
$\eta_W$ is the generic point of $W$.  In particular, $Z$ occurs with
coefficient $m_Z\geq1$.
Since the $H_i$ are globally generated, their mixed degree on every
component is non-negative.  Therefore
\begin{equation}\label{eq:component-degree}
 (H_1\cdots H_n\cdot Y)
 \geq m_Z(H_1\cdots H_n\cdot Z)
 \geq(H_1\cdots H_n\cdot Z).
\end{equation}

\smallskip
\noindent\emph{The mixed degree of the toric component.}
The $i$-th exponent polytope of $\tau$ is $P_i$.  General hyperplanes pull
back to general Laurent polynomials with Newton polytopes
$P_1,\ldots,P_n$.  Since their systems are basepoint-free and
$\dim(Z\setminus T)\leq n-1$, an incidence count shows that their common
zero set avoids the boundary.  The Bernstein--Kushnirenko theorem, or
equivalently the toric mixed-intersection formula \cite{Ber75,Ful93}, gives
\[
 (H_1\cdots H_n\cdot Z)=n!\vol(P_1,\ldots,P_n).
\]
Together with \eqref{eq:component-degree}, this proves
\eqref{eq:finite-level}.
\end{proof}

\begin{proposition}\label{prop:movable-upper}
Let $X$ be a smooth connected projective variety of dimension $n$, let $D_1,\ldots,D_n$ be big $\mathbb Q$-divisors, and let
$\nu\colon\mathbb C(X)^\times\to\mathbb Z^n_{\lex}$ be surjective and trivial on $\mathbb C^\times$. Then
\[
 n!\vol\Bigl(\Delta_\nu(D_1),\ldots,\Delta_\nu(D_n)\Bigr)\leq\langle D_1,\ldots,D_n\rangle.
\]
Here $\Delta_\nu(D)$ for a big $\mathbb Q$-divisor is defined by homogeneity after clearing denominators.
\end{proposition}

\begin{proof}
By homogeneity, assume that all $D_i$ are integral.  Let $x_\nu$ be the center of $\nu$ on $X$.  For each $i$, choose a non-zero rational section $\tau_i$ of $\mathcal O_X(D_i)$ which is regular and nowhere vanishing near $x_\nu$.
For every integer $k\geq1$, put
\[
 W_{i,k}\coloneqq \mathrm{H}^0\bigl(X,\mathcal O_X(kD_i)\bigr)\tau_i^{-k}
 \subset\mathbb C(X).
\]
Whenever $W_{i,k}\neq0$, put
\[
 P_{i,k}\coloneqq \Conv\nu(W_{i,k}^\times).
\]
Let $\Phi_k$ be the joint evaluation map associated with the $W_{i,k}$, as in \cref{lem:finite-level}:
\[
 \Phi_k=(\Phi_{1,k},\ldots,\Phi_{n,k})\colon X\dashrightarrow\prod_{i=1}^n\mathbb P(W_{i,k}^\vee).
\]
Under multiplication by $\tau_i^k$, its $i$-th component is the Kodaira map of $|kD_i|$.

The proof of \cite[Proposition~10.3.1]{Xiabook} supplies an integer $p_0>0$ and $s_0,s_1,\ldots,s_n\in W_{1,p_0}^\times$ with
\[
 \nu(s_j)-\nu(s_0)=e_j,\qquad 1\leq j\leq n.
\]
Choose a sufficiently divisible multiple $p=rp_0$ such that every map defined by $|pD_i|$ is birational.  Set
\[
 \widetilde s_0\coloneqq s_0^r,\qquad \widetilde s_j\coloneqq s_js_0^{r-1}\quad(1\leq j\leq n).
\]
These elements lie in $W_{1,p}^\times$ and satisfy $\nu(\widetilde s_j)-\nu(\widetilde s_0)=e_j$.  If $k=mp$, then
$\widetilde s_0^m,\widetilde s_j\widetilde s_0^{m-1}\in W_{1,k}^\times$, so
\[
 \operatorname{span}_{\mathbb Z}
 \{\nu(f)-\nu(g):f,g\in W_{i,k}^\times,\ 1\leq i\leq n\}
 =\mathbb Z^n.
\]
Moreover, the image of the multiplication map
\[
 \operatorname{Sym}^mW_{i,p}\longrightarrow W_{i,k}
\]
is a linear subsystem defining the $m$-th Veronese of the birational map associated with $|pD_i|$.  Hence $\Phi_{i,k}$, and therefore $\Phi_k$, is birational onto its image.

Let $Y_k$ be the closure of the image of $\Phi_k$, and let $H_{i,k}$ be the
restriction to $Y_k$ of the hyperplane class from the $i$-th factor.  Let
$\mu_k\colon X_k\to X$ be a common resolution of the base ideals of the
systems $|kD_i|$, with $X_k$ smooth, and write
\[
 \mu_k^*(kD_i)=M_{i,k}+F_{i,k},
 \qquad F_{i,k}\geq0,
\]
where $M_{i,k}$ is basepoint-free and $F_{i,k}$ is the effective fixed part.
The induced morphism $X_k\to Y_k$ is birational and pulls $H_{i,k}$ back to
$M_{i,k}$.  Hence the projection formula and
\cref{lem:finite-level} give
\begin{align*}
 n!\vol(P_{1,k},\ldots,P_{n,k})
 &\leq(H_{1,k}\cdots H_{n,k}\cdot Y_k)\\
 &=M_{1,k}\cdots M_{n,k}\\
 &\leq k^n\langle D_1,\ldots,D_n\rangle.
\end{align*}
The last inequality is \eqref{eq:movable-admissible}, since
$\mu_k^*D_i-k^{-1}M_{i,k}=k^{-1}F_{i,k}\geq0$.

Finally, \cite[Proposition~10.3.4]{Xiabook} gives
\[
 k^{-1}P_{i,k}=\Delta_{\nu,k}(D_i)
 \xrightarrow{\dHaus}\Delta_\nu(D_i).
\]
This holds in particular along $k\in p\mathbb Z_{>0}$.  Divide the preceding
inequality by $k^n$ and use continuity of mixed volume.
\end{proof}

\section{Uniformity in the valuation}

We first isolate a convex-geometric compactness statement.  It is important that neither the position nor the shape of the bodies is fixed.
Recall that $\mathcal{K}_n$ is the set of convex bodies in $\mathbb{R}^n$.

\begin{lemma}\label{lem:uniform-sandwich}
For $j\in\mathbb Z_{>0}$ and $1\leq i\leq n$, let $A_{i,j},C_{i,j},D_{i,j},B_{i,j}\in\Kconv$ satisfy
\[
 A_{i,j}\subseteq C_{i,j}\subseteq B_{i,j},
 \qquad
 A_{i,j}\subseteq D_{i,j}\subseteq B_{i,j}.
\]
Suppose that there are constants $v,M>0$, independent of $i$ and $j$, such that
\begin{align*}
 \vol(A_{i,j})&\geq v,\\
 \vol(B_{1,j}+\cdots+B_{n,j})&\leq M,\\
 \vol(B_{i,j})-\vol(A_{i,j})&\longrightarrow0
 \quad(1\leq i\leq n).
\end{align*}
Then
\[
 \left|\vol(C_{1,j},\ldots,C_{n,j})
       -\vol(D_{1,j},\ldots,D_{n,j})\right|\longrightarrow0.
\]
\end{lemma}
We refer to \cite{Sch14} for John's theorem and Blaschke selection theorem used in the proof below.
\begin{proof}
Translate each $i$-th quadruple by the same vector so that $0\in A_{i,j}$,
and set $S_j\coloneqq B_{1,j}+\cdots+B_{n,j}$.  These translations do not change
mixed volumes.  Every body is now contained in $S_j$, and
\[
 v\leq\vol(S_j)\leq M.
\]
Let $\mathbb B^n\subseteq\mathbb R^n$ be the closed Euclidean unit ball.
Choose an affine isomorphism
$\mathcal T_j(x)\coloneqq L_jx+b_j$, with
$L_j\in\operatorname{GL}_n(\mathbb R)$ and $b_j\in\mathbb R^n$, which
sends the John ellipsoid of $S_j$ to $\mathbb B^n$.  Writing
$\kappa_n\coloneqq \vol(\mathbb B^n)$, John's theorem gives
\[
 \mathbb B^n\subseteq \mathcal T_j(S_j)\subseteq n\mathbb B^n,
 \qquad
 \frac{\kappa_n}{M}\leq|\det L_j|
 \leq\frac{n^n\kappa_n}{v}.
\]
Thus the transformed bodies lie in one fixed compact set, the transformed
$A_{i,j}$ have volume bounded away from zero, and their transformed volume
gaps tend to zero.

Starting with any subsequence, Blaschke's selection theorem gives a further
subsequence along which all $4n$ transformed bodies converge.  At the limit,
\[
 \widehat A_i\subseteq\widehat C_i\subseteq\widehat B_i,
 \qquad
 \widehat A_i\subseteq\widehat D_i\subseteq\widehat B_i,
 \qquad \vol(\widehat A_i)=\vol(\widehat B_i)>0,
\]
where the hats denote the respective Hausdorff limits.  Hence
$\widehat A_i=\widehat C_i=\widehat D_i=\widehat B_i$.  Continuity of mixed
volume now makes the transformed mixed-volume difference tend to zero.
Finally,
\[
 \vol\Bigl(\mathcal T_j(K_1),\ldots,\mathcal T_j(K_n)\Bigr) =|\det L_j|\vol(K_1,\ldots,K_n)
\]
for arbitrary $K_1,\ldots,K_n\in\Kconv$; the affine translations are harmless.  The lower bound for
$|\det L_j|$ gives the same conclusion before transformation. Since this works for every subsequence, the full sequence converges.  
\end{proof}

The next result upgrades fixed-valuation continuity only at the level of mixed-volume values.  It does not assert valuation-uniform Hausdorff convergence of the bodies themselves.

\begin{proposition}\label{prop:uniform-dS}
For $1\leq i\leq n$, fix a big line bundle $L_i$, a smooth reference metric with curvature $\theta_i$, and potentials
$\varphi_{i,j},\varphi_i\in\PSH(X,\theta_i)_{>0}$ such that
\[
 \varphi_{i,j}\xrightarrow{\dS}\varphi_i.
\]
Then
\begin{equation}\label{eq:uniform-dS}
 \sup_\nu\Biggl|
 \vol\Bigl(\Delta_\nu(\theta_1,\varphi_{1,j}),\ldots,\Delta_\nu(\theta_n,\varphi_{n,j})\Bigr)-\vol\Bigl(\Delta_\nu(\theta_1,\varphi_1),\ldots,\Delta_\nu(\theta_n,\varphi_n)\Bigr)\Biggr|\longrightarrow 0,
\end{equation}
where the supremum is over all surjective valuations $\mathbb C(X)^\times\to\mathbb Z^n_{\lex}$ that are trivial on $\mathbb C^\times$.
\end{proposition}

\begin{proof}
Replace every $\varphi_{i,j}$ and $\varphi_i$ by its model envelope.  As recalled in \cref{sec:preliminaries}, this changes neither its $\dS$-class nor its partial Okounkov body.

Suppose that \eqref{eq:uniform-dS} fails.  Then there are a subsequence, still indexed by $j$, valuations $\nu_j$, and a constant $\delta>0$ for which the absolute difference in \eqref{eq:uniform-dS}, evaluated at $\nu_j$, is at least $\delta$.

Apply \cref{lem:potential-sandwich} successively for
$i=1,\ldots,n$, at each step restricting the subsequence obtained at the
preceding step.  After this finite diagonal extraction, there are increasing
model potentials $\eta_{i,j}$ and decreasing model potentials $\psi_{i,j}$
on one common subsequence such that
\begin{equation}\label{eq:potential-sandwich}
 \eta_{i,j}\leq\varphi_{i,j}\leq\psi_{i,j},
 \qquad
 \eta_{i,j}\leq\varphi_i\leq\psi_{i,j},
 \qquad
 \eta_{i,j}\xrightarrow{\dS}\varphi_i,
 \quad
 \psi_{i,j}\xrightarrow{\dS}\varphi_i.
\end{equation}

By the $\dS$-continuity of raw non-pluripolar mass
\cite[Corollary~6.2.5]{Xiabook}, the diagonal masses of
$\eta_{i,j}$ and $\psi_{i,j}$ converge to the positive raw mass of
$\varphi_i$.  We may therefore discard one finite initial segment so that
all partial Okounkov bodies occurring below lie on the positive-volume
locus.

For every $i,j$, put
\begin{align*}
 A_{i,j}&\coloneqq \Delta_{\nu_j}(\theta_i,\eta_{i,j}),&
 B_{i,j}&\coloneqq \Delta_{\nu_j}(\theta_i,\psi_{i,j}),\\
 C_{i,j}&\coloneqq \Delta_{\nu_j}(\theta_i,\varphi_{i,j}),&
 D_{i,j}&\coloneqq \Delta_{\nu_j}(\theta_i,\varphi_i).
\end{align*}
By \eqref{eq:body-monotonicity} and \eqref{eq:potential-sandwich},
\[
 A_{i,j}\subseteq C_{i,j}\subseteq B_{i,j},
 \qquad
 A_{i,j}\subseteq D_{i,j}\subseteq B_{i,j}.
\]
The numerical bounds in \cref{lem:uniform-sandwich} do not depend on $\nu_j$. Indeed, \eqref{eq:partial-body-volume} and the $\dS$-continuity of the volume \cite[Corollary~7.3.1]{Xiabook} give
\begin{align*}
 \vol(B_{i,j})-\vol(A_{i,j})
 &=\frac1{n!}\Bigl(\vol(\theta_i+\ddc\psi_{i,j})-\vol(\theta_i+\ddc \eta_{i,j})\Bigr)\longrightarrow 0,\\
 \vol(A_{i,j})
 &=\frac{1}{n!}\vol(\theta_i+\ddc\eta_{i,j})
 \longrightarrow\frac{1}{n!}\vol(\theta_i+\ddc\varphi_i)>0.
\end{align*}
Furthermore, \eqref{eq:body-Minkowski} and \eqref{eq:body-monotonicity} yield
\[
 B_{1,j}+\cdots+B_{n,j}
 \subseteq
 \Delta_{\nu_j}\left(\sum_i\theta_i+\ddc \psi_{i,j}\right)
 \subseteq\Delta_{\nu_j}(L_1\otimes\cdots\otimes L_n).
\]
By \eqref{eq:ordinary-volume}, the last body has Euclidean volume $\vol(L_1\otimes\cdots\otimes L_n)/n!$, independently of $j$ and $\nu_j$.

All assumptions of \cref{lem:uniform-sandwich} are therefore satisfied.  It forces the mixed-volume difference at $\nu_j$ to tend to zero, contradicting its lower bound $\delta$.
\end{proof}

\section{Proof of the mixed-volume formula}

\begin{proposition}\label{prop:analytic-case}
The conclusion of \cref{thm:main} holds if each $h_i$ has analytic singularities and every $c_1(L_i,h_i)$ is a K\"ahler current.
\end{proposition}

\begin{proof}
Take a projective log resolution $\pi\colon Y\to X$, with $Y$ smooth and projective, for which
\[
 \pi^*c_1(L_i,h_i)=[D_i]+R_i,
\]
where $D_i$ is an effective $\mathbb Q$-divisor and $R_i$ is a closed positive current with locally bounded potentials; see
\cite[Definition~1.6.2 and Theorem~1.6.1]{Xiabook}.  Put
\[
 \alpha_i\coloneqq\{R_i\}
 =\pi^*c_1(L_i)-\{D_i\}\in N^1(Y)_{\mathbb Q}.
\]
The class $\alpha_i$ is nef by \cite[Lemma~1.6.1]{Xiabook}.  Choose a
nef $\mathbb Q$-Cartier divisor in this numerical class and denote it again
by $\alpha_i$.  This is harmless because ordinary Okounkov bodies depend
only on numerical classes for the present valuations
\cite[Proposition~10.3.2]{Xiabook}.  Then
\begin{equation}\label{eq:mass-on-resolution}
 \int_Xc_1(L_1,h_1)\wedge\cdots\wedge c_1(L_n,h_n)=\int_Y R_1\wedge\cdots\wedge R_n=\alpha_1\cap \cdots \cap \alpha_n.
\end{equation}
Here we used \cite[Proposition~2.4.2 and Proposition~2.2.1(1), (3), and (4)]{Xiabook}. Note that
\[
 \alpha_i^n=\int_X c_1(L_i,h_i)^n>0,
\]
so the nef divisor $\alpha_i$ is big.

For a $\mathbb Q$-Cartier divisor $D$ and a valuation $\nu$ centered at $y_\nu\in Y$, choose $m>0$ such that $mD$ is Cartier and let $g$ be a local equation of $mD$ at $y_\nu$. Set
\[
 \nu(D)\coloneqq m^{-1}\nu(g)\in\mathbb Q^n.
\]
For every valuation $\nu$, the log-resolution formula \cite[Remark~10.3.1, Equation~(10.22)]{Xiabook} reads
\begin{equation}\label{eq:resolution-body}
 \Delta_\nu(L_i,h_i)=\Delta_\nu(\alpha_i)+\nu(D_i).
\end{equation}
Translation invariance, \cref{prop:movable-upper}, and nefness give
\begin{equation}\label{eq:analytic-upper}
 n!\vol\Bigl(\Delta_\nu(L_1,h_1),\ldots,\Delta_\nu(L_n,h_n)\Bigr) \leq\langle \alpha_1,\ldots,\alpha_n\rangle=\alpha_1\cap \cdots \cap \alpha_n.
\end{equation}

It remains to prove the opposite inequality after taking the supremum. Fix an ample $\mathbb Q$-divisor $H$ on $Y$. Choose one
$c\in\mathbb Q_{>0}$ so large that every $c\alpha_i-H$ is big, and choose effective $\mathbb Q$-divisors
\[
 E_i\sim_{\mathbb Q}c\alpha_i-H.
\]
For $\epsilon\in\mathbb Q_{>0}$ set
\[
 \beta_i\coloneqq \alpha_i+\epsilon H.
\]
The $\mathbb Q$-divisor $\beta_i$ is ample and
\begin{equation}\label{eq:Kodaira-decomposition}
 \frac{1}{1+c\epsilon}\beta_i
 +\frac{\epsilon}{1+c\epsilon}E_i
 \sim_{\mathbb Q}\alpha_i.
\end{equation}
By \cite[Theorem~1.1]{Wil25}, there is an admissible flag
\[
 Y_\bullet\colon Y=Y_0\supset Y_1\supset\cdots\supset Y_n=\{y\},
\]
where each $Y_r$ is irreducible of codimension $r$ in $Y$ and is nonsingular at $y$.  Let $\nu_\epsilon$ be the valuation defined by this
flag.  It
satisfies
\begin{equation}\label{eq:Wilms}
 n!\vol\Bigl(\Delta_{\nu_\epsilon}(\beta_1),\ldots,\Delta_{\nu_\epsilon}(\beta_n)\Bigr)=\beta_1\cap \cdots \cap \beta_n.
\end{equation}
The flag valuation is surjective and trivial on $\mathbb C^\times$; since $\mathbb C(Y)=\mathbb C(X)$, it is among the valuations in \cref{thm:main}.  After clearing denominators, multiplication by the canonical section of $E_i$ gives
\[
 \Delta_{\nu_\epsilon}(\beta_i)
 +\epsilon\nu_\epsilon(E_i)
 \subseteq
 \Delta_{\nu_\epsilon}(\beta_i+\epsilon E_i).
\]
By homogeneity, numerical invariance
\cite[Proposition~10.3.2]{Xiabook}, and
\eqref{eq:Kodaira-decomposition}, this becomes
\[
 \frac{1}{1+c\epsilon}\Delta_{\nu_\epsilon}(\beta_i)+\frac{\epsilon}{1+c\epsilon}\nu_\epsilon(E_i)\subseteq\Delta_{\nu_\epsilon}(\alpha_i).
\]
Using \eqref{eq:Wilms}, we find
\begin{equation}\label{eq:analytic-lower-approx}
 \sup_\nu\vol\Bigl(\Delta_\nu(\alpha_1),\ldots,\Delta_\nu(\alpha_n)\Bigr)
 \geq\frac{1}{n!(1+c\epsilon)^n}\beta_1\cdots \beta_n.
\end{equation}
Letting $\epsilon\to 0+$ in \eqref{eq:analytic-lower-approx}, and then using the translations in \eqref{eq:resolution-body}, proves the reverse of \eqref{eq:analytic-upper}. Together with \eqref{eq:mass-on-resolution}, this proves the proposition.
\end{proof}

\begin{proof}[Proof of \cref{thm:main}]
Choose smooth Hermitian metrics $h_i^0$ on $L_i$ and write
\[
 \theta_i\coloneqq c_1\left(L_i,h_i^0\right),\qquad
 h_i=h_i^0\mathrm e^{-\varphi_i}.
\]
Set
\[
 \widetilde\varphi_i\coloneqq
 P_{\theta_i}[\varphi_i]_{\mathcal I}.
\]
By \cite[Proposition~3.2.6]{Xiabook}, $\varphi_i$ and
$\widetilde\varphi_i$ are $\mathcal I$-equivalent, so their partial
Okounkov bodies agree for every $\nu$.  By the definition of Cao's mixed
volume, the left-hand side of \eqref{eq:main} is the mixed non-pluripolar
mass of the $\widetilde\varphi_i$.  Moreover,
\[
 \int_X\theta_{i,\widetilde\varphi_i}^n
 =\vol(\theta_i+\ddc\varphi_i)>0,
\]
and each $\widetilde\varphi_i$ is an $\mathcal I$-good model potential by
\cite[Example~7.1.2]{Xiabook}.  We may therefore replace the metrics by
$h_i^0\mathrm e^{-\widetilde\varphi_i}$ and relabel
$\widetilde\varphi_i$ as $\varphi_i$; this changes neither side of
\eqref{eq:main} and places every potential in
$\PSH(X,\theta_i)_{>0}$.

\Cite[Theorem~7.1.1]{Xiabook} gives potentials $\psi_{i,j}$ with analytic singularities such that
\[
 \psi_{i,j}\xrightarrow{\dS}\varphi_i,
 \qquad \theta_{i,\psi_{i,j}}\ \text{is a K\"ahler current}.
\]
By \cref{prop:analytic-case}, for every $j$,
\begin{equation}\label{eq:approximating-equality}
 \int_X\theta_{1,\psi_{1,j}}\wedge\cdots\wedge\theta_{n,\psi_{n,j}}=n!\sup_\nu\vol\Bigl(
  \Delta_\nu(\theta_1,\psi_{1,j}),\ldots,
  \Delta_\nu(\theta_n,\psi_{n,j})\Bigr).
\end{equation}
By \cite[Theorem~6.2.1]{Xiabook}, the left-hand side converges to
\[
 \int_Xc_1(L_1,h_1)\wedge\cdots\wedge c_1(L_n,h_n).
\]
By \cref{prop:uniform-dS}, the right-hand side converges to
\[
 n!\sup_\nu\vol\Bigl(\Delta_\nu(L_1,h_1),\ldots,\Delta_\nu(L_n,h_n)\Bigr).
\]
Indeed, if $f_j(\nu)$ and $f(\nu)$ denote the corresponding mixed-volume
values, then
\[
 \left|\sup_\nu f_j(\nu)-\sup_\nu f(\nu)\right|
 \leq\sup_\nu|f_j(\nu)-f(\nu)|\longrightarrow0
\]
by \cref{eq:uniform-dS}.
Taking limits in \eqref{eq:approximating-equality} proves \eqref{eq:main}.
\end{proof}

\begin{proof}[Proof of \cref{thm:main1}]
For each $i$, choose a psh metric $h_i$ with minimal singularities on
$L_i$.  The partial body is contained in the ordinary body,
\[
 \Delta_\nu(L_i,h_i)\subseteq\Delta_\nu(L_i).
\]
The diagonal formulas \eqref{eq:partial-diagonal-volume} and
\eqref{eq:diagonal-volume}, together with
$\vol\ddc h_i=\vol(L_i)$, show that these two compact convex sets have the
same positive Euclidean volume.  Strict volume monotonicity for
full-dimensional convex bodies therefore gives
\[
 \Delta_\nu(L_i,h_i)=\Delta_\nu(L_i)
\]
for every valuation $\nu$ under consideration.  On the other hand,
\cite[Definition~1.17]{BEGZ10} identifies the mixed non-pluripolar product
of currents with minimal singularities with the positive intersection
product; in the projective numerical setting this is the movable product
of \cite[Propositions~2.12--2.13]{BFJ09}.  Applying \cref{thm:main} to the
metrics $h_i$ now turns its left-hand side into
$\langle L_1,\ldots,L_n\rangle$ and its partial bodies into the ordinary
bodies.  The resulting equality is \eqref{eq:movable-main}.
\end{proof}

\printbibliography
Mingchen Xia, \textsc{Institute of Geometry and Physics, USTC}\par\nopagebreak
  \textit{Email address:} \texttt{xiamingchen2008@gmail.com}
\end{document}
